% ---------- 4. Searle and understanding ----------
\begin{frame}{Searle's Chinese Room (1980)}
\begin{columns}[T]
\begin{column}{0.36\textwidth}
\small
\textbf{Thought experiment.} You, who understand no Chinese, sit in a closed room with a big rule book (in English) and baskets of Chinese symbols.

Slips with Chinese questions come in. You follow the rules, match shapes, and pass out other slips.

To the people outside, your answers are \textbf{fluent Chinese}. Inside, you understand \textbf{nothing}.

\vspace{1mm}
\textbf{Searle's point:} a program only manipulates symbols by their \emph{form} (syntax). That cannot, by itself, give \emph{meaning} (semantics).
\end{column}
\begin{column}{0.61\textwidth}
\begin{sfig}[0.7]
\begin{tikzpicture}[scale=1.0,every node/.append style={font=\small}]
  \draw[rounded corners=6pt,thick,draw=tdInk!70,fill=tdPaper!60!tdLine!30] (2.2,-1.9) rectangle (8.2,1.9);
  \node[font=\scriptsize\bfseries,anchor=north west] at (2.3,1.85) {the room};
  \node[card=tdBlue,minimum width=1.6cm] (p) at (4.0,0) {\textbf{you}\\{\scriptsize no Chinese}};
  \node[card=tdAmber,minimum width=1.7cm] (r) at (6.6,0.7) {\textbf{rule book}\\\shortstack{\scriptsize if shape X,\\ send shape Y}};
  \node[card=tdGreen,minimum width=1.7cm] (s) at (6.6,-0.9) {\textbf{symbols}\\{\scriptsize baskets}};
  \draw[arr] (p) -- (r); \draw[arr] (p) -- (s);
  \node[soft,minimum width=1.5cm] (in) at (-0.4,0.8) {\scriptsize\ttfamily ni hao ma?};
  \node[soft,minimum width=1.5cm] (out) at (10.5,-0.8) {\scriptsize\ttfamily wo hen hao};
  \draw[arr] (in) -- (2.2,0.4); \draw[arr] (8.2,-0.4) -- (out);
  \node[lbl] at (-0.4,1.6) {question in}; \node[lbl] at (10.5,0.0) {answer out};
\end{tikzpicture}
\end{sfig}
\end{column}
\end{columns}
\src{Searle, ``Minds, Brains, and Programs'', Behavioral and Brain Sciences 3(3), 417--457, 1980.}
\end{frame}

\begin{frame}{The argument, step by step}
\begin{columns}[T]
\begin{column}{0.55\textwidth}
\small
\begin{enumerate}\setlength\itemsep{4pt}
\item A computer program is purely \textbf{syntactic}: it manipulates symbols by their form.
\item Minds have \textbf{semantics}: thoughts are \emph{about} things.
\item Syntax is not sufficient for semantics (the room shows this).
\item \textbf{So} running a program is not sufficient for having a mind.
\end{enumerate}
\vspace{1mm}
Searle's targets: \textbf{strong AI} (the right program \emph{is} a mind) versus \textbf{weak AI} (programs are useful tools for studying or simulating minds). He accepts weak AI.

{\footnotesize He does \emph{not} say machines cannot think: brains are physical machines that cause minds. He says the \emph{causal powers of brains} matter, not just the program.}
\end{column}
\begin{column}{0.42\textwidth}
\begin{sfig}[0.6]
\begin{tikzpicture}
  \node[card=tdBlue,minimum width=4cm] (a) at (0,2.2) {\textbf{Programs}\\{\scriptsize syntax only}};
  \node[card=tdGreen,minimum width=4cm] (b) at (0,0.2) {\textbf{Minds}\\{\scriptsize have meaning (semantics)}};
  \node[card=tdRed,minimum width=4cm] (c) at (0,-2.0) {\textbf{Conclusion}\\\shortstack{\scriptsize programs alone do\\ not suffice for minds}};
  \draw[arr] (a) -- node[right,lbl] {syntax $\neq$ semantics} (c);
  \draw[arr] (b) -- (c);
\end{tikzpicture}
\end{sfig}
\end{column}
\end{columns}
\end{frame}

\begin{frame}[fragile]{Demo 2: a Chinese Room in a dozen lines}
\begin{columns}[T]
\begin{column}{0.5\textwidth}
\pycodes{chinese_room_room}
\vspace{1mm}
{\footnotesize The rule book is a Python dictionary: 4 rules, written in romanised Chinese.}
\end{column}
\begin{column}{0.47\textwidth}
\pyouts{chinese_room_run}
\vspace{1mm}
\pyouts{chinese_room_stats}
\vspace{1mm}
{\footnotesize On the four known slips the room \emph{looks} competent. On a new question it stalls with ``please say that again''. A real room would need an astronomically large book, which is exactly what the replies below argue about.}
\end{column}
\end{columns}
\end{frame}

\begin{frame}{Four classic replies, and Searle's answers}
\begin{sfig}[0.66]
\begin{tikzpicture}
  \node[card=tdBlue,text width=6.4cm,minimum height=2.8cm] at (0,0) {\textbf{Systems reply}\\{\footnotesize The \emph{whole system} (you + book + room) understands, not the man alone.}\\[3pt]{\scriptsize\color{tdRed!70!black}Searle: let the man memorise everything; still no understanding.}};
  \node[card=tdGreen,text width=6.4cm,minimum height=2.8cm] at (7.2,0) {\textbf{Robot reply}\\{\footnotesize Put the program in a robot with cameras and limbs: symbols get \emph{grounded} in the world.}\\[3pt]{\scriptsize\color{tdRed!70!black}Searle: camera input is more symbols; the man still just shuffles them.}};
  \node[card=tdAmber,text width=6.4cm,minimum height=2.8cm] at (0,-3.3) {\textbf{Brain simulator reply}\\{\footnotesize Simulate every neuron of a Chinese speaker's brain.}\\[3pt]{\scriptsize\color{tdRed!70!black}Searle: simulating a storm does not make you wet; simulation is not duplication.}};
  \node[card=tdPurple,text width=6.4cm,minimum height=2.8cm] at (7.2,-3.3) {\textbf{Other minds reply}\\{\footnotesize We attribute understanding to people only from behaviour. Same standard for machines?}\\[3pt]{\scriptsize\color{tdRed!70!black}Searle: the issue is what understanding \emph{is}, not how we find out about it.}};
\end{tikzpicture}
\end{sfig}
\discuss{Which reply do you find strongest? Which of Searle's answers feels like a dodge?}
\src{Searle (1980) and the open peer commentary printed with it; Stanford Encyclopedia of Philosophy, ``The Chinese Room Argument''.}
\end{frame}

\begin{frame}{Do large language models understand?}
\begin{sfig}[0.55]
\begin{tikzpicture}
  \node[card=tdRed,text width=7.0cm,minimum height=3.7cm,anchor=north east] at (-0.2,0) {\textbf{``Stochastic parrots''}\\[2pt]\footnotesize Bender et al.\ (2021): LMs stitch together text by probability, without reference to meaning. Bender and Koller (2020): form alone cannot teach meaning (the ``octopus test'').\\[2pt]{\scriptsize\color{tdMuted} Warns of bias, cost, and people mistaking fluency for a mind.}};
  \node[card=tdGreen,text width=7.0cm,minimum height=3.7cm,anchor=north west] at (0.2,0) {\textbf{``Something more is going on''}\\[2pt]\footnotesize LLMs solve novel problems, use abstractions and show internal structure: world-model-like features (board-game and spatial probes), planning in circuits. Prediction at scale may force useful internal models.\\[2pt]{\scriptsize\color{tdMuted} Evidence is partial and contested.}};
  \node[pill=tdAmber] at (0,-4.2) {middle positions exist: some competence, little grounding, unclear what ``understanding'' means};
\end{tikzpicture}
\end{sfig}
\discuss{What observable evidence would convince you that a model understands? What would convince you it does not?}
\src{Bender and Koller, ACL 2020; Bender, Gebru, McMillan-Major, Shmitchell, ``On the Dangers of Stochastic Parrots'', FAccT 2021; Li et al., ``Emergent World Representations'', ICLR 2023.}
\end{frame}

\begin{frame}[fragile]{Demo 3: a tiny ``parrot''}
\begin{columns}[T]
\begin{column}{0.52\textwidth}
\pycodes{parrot_parrot}
\end{column}
\begin{column}{0.45\textwidth}
\pyouts{parrot_run}
\vspace{1mm}
{\footnotesize
Each word is chosen using \textbf{only the previous word}. The text is locally plausible and globally empty.

\textbf{Careful:} an LLM is \emph{not} a bigram table. It conditions on thousands of earlier tokens and has billions of learned parameters. The demo shows what ``predicting the next word'' alone does \emph{not} prove: that fluency needs no understanding at tiny scale. Whether that stays true at scale is the debate.}
\end{column}
\end{columns}
\end{frame}

\begin{frame}{Emergent abilities: real or a mirage?}
\begin{columns}[T]
\begin{column}{0.5\textwidth}
\small
\textbf{Claim (Wei et al.\ 2022).} Some abilities (multi-step arithmetic, certain reasoning tasks) are \textbf{absent in small models and appear abruptly} past a scale threshold.

\textbf{Challenge (Schaeffer et al.\ 2023).} Much of the ``jump'' comes from the \textbf{metric}: all-or-nothing scores (exact match) hide smooth underlying progress. Use a continuous metric and the cliff often becomes a slope.

\vspace{1mm}
Take-away: be careful with stories about sudden ``sparks''. Ask what was measured.
\end{column}
\begin{column}{0.47\textwidth}
\begin{sfig}[0.6]
\begin{tikzpicture}
\begin{axis}[width=7.4cm,height=5.2cm,xlabel={model scale (log)},ylabel={task score},ymin=-0.05,ymax=1.1,xmin=0,xmax=10,xtick=\empty,ytick={0,0.5,1},legend pos=north west]
\addplot[very thick,tdRed,domain=0:10,samples=60] {1/(1+exp(-(x-6.5)*3.0))};
\addplot[very thick,tdGreen,domain=0:10,samples=60] {0.05+0.9*x/10};
\legend{exact-match metric,continuous metric}
\end{axis}
\end{tikzpicture}
\end{sfig}
{\scriptsize\illus: the same hypothetical improvement, scored two ways.}
\end{column}
\end{columns}
\src{Wei et al., ``Emergent Abilities of Large Language Models'', TMLR 2022; Schaeffer, Miranda, Koyejo, ``Are Emergent Abilities of Large Language Models a Mirage?'', NeurIPS 2023.}
\end{frame}

\begin{frame}{Looking inside: mechanistic interpretability}
\begin{columns}[T]
\begin{column}{0.52\textwidth}
\small
Instead of arguing from outside, \textbf{open the box}: find the features and circuits a network uses.

Anthropic's March 2025 \emph{circuit tracing} study of Claude 3.5 Haiku reported, for example:
\begin{itemize}\setlength\itemsep{3pt}
\item when writing a rhyming couplet, the model \textbf{chooses the rhyme word before writing the line}; editing that internal ``plan'' changes the ending;
\item shared, language-independent concepts across languages;
\item in one hard maths case, a plausible-sounding explanation that \textbf{did not match} the computation actually done.
\end{itemize}
\end{column}
\begin{column}{0.45\textwidth}
\begin{sfig}[0.3]\begin{tikzpicture}
  \node[sk=tdBlue,text width=2.0cm] (a) at (0,0) {line 1 of a couplet ends};
  \node[sk=tdAmber,text width=2.2cm] (b) at (3.4,0) {internal feature: rhyme word chosen};
  \node[sk=tdGreen,text width=2.0cm] (c) at (6.8,0) {line 2 written toward it};
  \draw[arr] (a) -- (b); \draw[arr] (b) -- (c);
  \node[font=\tiny,text=tdRed!70!black] at (3.4,-1.0) {suppress ``rabbit'': ending changes};
\end{tikzpicture}\end{sfig}
\begin{block}{Honest limits (from the authors)}
{\small
\begin{itemize}\setlength\itemsep{2pt}
\item only a \textbf{fraction} of the computation;
\item hours of effort per \textbf{short prompt};
\item possible artefacts of the method.
\end{itemize}}
\end{block}
{\footnotesize Planning and internal representations count \emph{against} the ``pure parrot'' picture, but say nothing about \textbf{consciousness}.}
\end{column}
\end{columns}
\src{Anthropic, ``Tracing the thoughts of a large language model'' (27 Mar 2025) and ``On the Biology of a Large Language Model'', transformer-circuits.pub (2025); open-source tooling announced 29 May 2025.}
\end{frame}
